\chapter{2023年上海卷（秋）}

\section{填空题}

\begin{problem}
设 \(x\in\R\)，不等式 \(\left|x-2\right|<1\) 的解集为\fillinblank{}.
\end{problem}

\begin{problem}
若向量 \(\bs a=\left(-2,3\right)\)，\(\bs b=\left(1,2\right)\)，则 \(\bs a\cdot\bs b=\fillinblank{}\).
\end{problem}

\begin{problem}
若数列 \(\{a_n\}\) 是首项为 \(3\)，公比为 \(2\) 的等比数列，其前 \(n\) 项和为 \(S_n\)，则 \(S_6=\fillinblank{}\).
\end{problem}

\begin{problem}
若 \(\tan A=3\)，则 \(\tan 2A=\fillinblank{}\).
\end{problem}

\begin{problem}
函数 \(y=\begin{cases}
2^x,&x>0,\\
1,&x\le 0,
\end{cases}\) 的值域为\fillinblank{}.
\end{problem}

\begin{problem}
若复数 \(z=1+\i\)（\(\i\) 为虚数单位），则 \(\left|1-\i z\right|=\fillinblank{}\).
\end{problem}

\begin{problem}
设 \(m\in\R\). 若圆 \(x^2+y^2-4y-m=0\) 的面积为 \(\pi\)，则 \(m=\fillinblank{}\).
\end{problem}

\begin{problem}
在 \(\triangle ABC\) 中，角 \(A\)，\(B\)，\(C\) 所对的边长分别为 \(a\)，\(b\)，\(c\). 若 \(a=4\)，\(b=5\)，\(c=6\)，则 \(\sin A=\fillinblank{}\).
\end{problem}

\begin{problem}
国内生产总值（GDP）是衡量地区经济状况的最佳指标. 根据统计数据显示，某市在 2020 年间经济高质量增长，GDP 稳步增长，第一季度和第四季度的 GDP 分别为 \(232\) 亿元和 \(241\) 亿元，且四个季度 GDP 逐季度增长. 若这四个季度 GDP 的中位数与平均数相等，则该市 2020 年的 GDP 总额为\fillinblank{} 亿元.
\end{problem}

\begin{problem}
已知对任意给定的实数 \(x\)，都有
\(
\left(1+2023x\right)^{100}+\left(2023-x\right)^{100}=a_0+a_1x+a_2x^2+\cdots+a_{100}x^{100},
\)
其中 \(a_0\)，\(a_1\)，\(\cdots\)，\(a_{100}\in\R\). 若 \(a_k<0\)，其中 \(k\in\{0,1,2,\cdots,100\}\)，则正整数 \(k\) 的最大值为\fillinblank{}.
\end{problem}

\begin{problem}
某公园欲建设一段斜坡，斜坡面是一平面，坡面与水平地面所成角为 \(\theta\)，坡顶距离水平地面的垂直高度为 \(4\) 米，游客从坡底沿着斜坡面向上直行，每走 \(1\) 米消耗的体力为 \(1.025-\cos\theta\)，当 \(\theta=\fillinblank{}\) 时，游客走斜坡所消耗的总体力最小.
\end{problem}

\begin{problem}
空间中有三个定点 A，B，C，且 \(AB=BC=CA=1\). 若在空间中任取 \(2\) 个不同的点，使它们与 A，B，C 恰好成为一个正四棱锥的五个顶点，则不同的取法共有\fillinblank{} 种.
\end{problem}

\section{单选题}

\begin{problem}
已知集合 \(P=\{1,2\}\)，\(Q=\{2,3\}\). 若 \(M=\{x\mid x\in P,\ x\notin Q\}\)，则 \(M=\quad\).

\choices
    {\(\{1\}\)}
    {\(\{2\}\)}
    {\(\{3\}\)}
    {\(\{1,2,3\}\)}
\end{problem}

\begin{problem}
为了解某校高中男生身高与体重的关系，随机选取 \(50\) 名男生的身高与体重的数据，绘制散点图，如图所示，下列说法正确的是（\quad）.

\bitmapfigure[max width=0.68\linewidth]{img/2023/shanghai/q14_fig1.png}

\choices
    {身高越高，体重越重}
    {身高越高，体重越轻}
    {身高与体重之间呈正相关}
    {身高与体重之间呈负相关}
\end{problem}

\begin{problem}
已知 \(a>0\)，函数 \(y=\sin x\) 在区间 \(\left[a,2a\right]\) 上的最小值为 \(s\)，在区间 \(\left[2a,3a\right]\) 上的最小值为 \(t\)，当 \(a\) 变化时，下列一定不成立的是（\quad）.

\choices
    {\(s>0,t>0\)}
    {\(s>0\)，\(t<0\)}
    {\(s<0\)，\(t<0\)}
    {\(s<0\)，\(t>0\)}
\end{problem}

\begin{problem}
对于曲线 \(C\)，若存在点 \(M\)，使得对任意的点 \(P\in C\)，都存在 \(Q\in C\)，且 \(\left|PM\right|\cdot\left|QM\right|=1\)，则称曲线 \(C\) 为“自相关曲线”. 有以下结论：
\begin{circlelist}
\item 任何椭圆都是“自相关曲线”；
\item 存在双曲线是“自相关曲线”. 关于上述两个结论，说法正确的是（\quad）.
\end{circlelist}

\choices
    {\circled{1} 成立，\circled{2} 成立}
    {\circled{1} 成立，\circled{2} 不成立}
    {\circled{1} 不成立，\circled{2} 成立}
    {\circled{1} 不成立，\circled{2} 不成立}
\end{problem}

\section{解答题}

\begin{problem}
2023 年 6 月 7 日，21 世纪汽车博览会在上海举行. 已知某汽车模型公司共有 \(25\) 个汽车模型，其外观和内饰的颜色分布如下表所示：

\begin{center}
\begin{tabular}{c|cc}
\hline
& 红色外观 & 蓝色外观\\
\hline
米色内饰 & \(12\) & \(8\)\\
棕色内饰 & \(2\) & \(3\)\\
\hline
\end{tabular}
\end{center}

\begin{enumerate}
\item 小明从这些汽车模型中随机取出一个，记事件 A 为“小明取到红色外观的汽车模型”，事件 B 为“小明取到棕色内饰的汽车模型”，求 \(P(B)\) 和 \(P(B\mid A)\)，并据此判断事件 A 和事件 B 是否独立；
\item 该公司举行抽奖活动，规定在一次抽奖中，每人可以一次性从这些汽车模型中随机取出两个汽车模型，并给出以下假设：
\begin{enumerate}
\item 取出的两个汽车模型会出现三种结果，即外观和内饰均为同色，外观和内饰均不为同色，以及仅外观或仅内饰为同色；
\item 按出现结果的可能性大小，概率越小奖项越高；
\item 抽奖活动的奖金金额为：一等奖 \(600\) 元，二等奖 \(300\) 元，三等奖 \(150\) 元.
\end{enumerate}
请分析各奖项对应的结果，设 \(X\) 为奖金额，写出 \(X\) 的分布并求其数学期望.
\end{enumerate}
\end{problem}

\begin{problem}
在四棱柱 \(ABCD-A_1B_1C_1D_1\) 中，\(AB\parallel DC\)，\(AB\perp AD\)，\(AB=2\)，\(AD=3\)，\(DC=4\).

\bitmapfigure[max width=0.36\linewidth]{img/2023/shanghai/q17_fig1.png}
\end{problem}

\begin{problem}
设函数 \(f(x)=\frac{x^2+(3a+1)x+c}{x+a}\)，\(a\)，\(c\in\R\)
\begin{enumerate}
\item 当 \(a=0\) 时，是否存在实数 \(c\)，使得 \(y=f(x)\) 为奇函数，请说明理由；
\item 若函数 \(y=f(x)\) 的图像过点 \((1,3)\)，且其与 \(x\) 轴的负半轴有两个不同的交点，求 \(c\) 的值及 \(a\) 的取值范围.
\end{enumerate}
\end{problem}

\begin{problem}
已知抛物线 \(\Gamma:y^2=4x\)，第一象限内点 \(A\) 在 \(\Gamma\) 上，\(A\) 的纵坐标是 \(a\).
\begin{enumerate}
\item 若 \(A\) 到准线的距离为 \(3\)，求 \(a\)；
\item 若 \(a=4\)，\(B\) 在 \(x\) 轴上，且线段 \(AB\) 的中点在 \(\Gamma\) 上，求点 \(B\) 的坐标及坐标原点 \(O\) 到直线 \(AB\) 的距离；
\item 直线 \(l:x=-3\)，令 \(P\) 是第一象限 \(\Gamma\) 上异于 A 的一点，直线 \(AP\) 交 \(l\) 于 \(Q\)，\(H\) 是 \(P\) 在 \(l\) 上的投影，若点 A 满足“对于任意 \(P\) 都有 \(\left|HQ\right|>4\)”，求 \(a\) 的取值范围.
\end{enumerate}
\end{problem}

\begin{problem}
已知 \(f(x)=\ln x\). 曲线 \(y=f(x)\) 在点 \(\left(a_1,f(a_1)\right)\) 处的切线交 \(y\) 轴于点 \((0,a_2)\)，且 \(a_2>0\)；曲线 \(y=f(x)\) 在点 \(\left(a_2,f(a_2)\right)\) 处的切线交 \(y\) 轴于点 \((0,a_3)\)；以此类推，直至 \(a_m\le 0\) 时停止，得到数列 \(\{a_n\}\).
\begin{enumerate}
\item 证明：当正整数 \(n\ge 2\) 时，\(a_n=\ln a_{n-1}-1\)；
\item 当正整数 \(n\ge 2\) 时，试比较 \(a_n\) 与 \(a_{n-1}-2\) 的大小；
\item 若正整数 \(k\ge 3\)，是否存在 \(k\) 使得 \(a_1\)，\(a_2\)，\(\cdots\)，\(a_k\) 依次成等差数列？若存在，求出 \(k\) 的所有取值；若不存在，试说明理由.
\end{enumerate}
\end{problem}

